\chapter{Relations}
\label{ch:relations}

Sets don't care about order: $\set{a, b} = \set{b, a}$. To compare, combine or map
objects we need a pair in which one element comes \emph{first}. Then we use pairs
to say when objects are \emph{related}, and from that we get the two most
important kinds of relations: \emph{equivalences} and \emph{orders}.

\section{Ordered pairs}

\begin{definition}[Ordered pair, Kuratowski]
\label{def:pair}
\[
  (a, b) \coloneqq \set{\set{a}, \set{a, b}}.
\]
$a$ is the \emph{first component} and $b$ the \emph{second component}.
\register{relations}{$(a, b)$}{ordered pair with first component $a$ and second
component $b$}
\end{definition}

This exists by \cref{ax:pairing}. The definition is only useful because of the
following property, which is all we will ever use about pairs.

\begin{theorem}[Characteristic property of pairs]
\label{thm:pair}
\[
  (a, b) = (c, d) \;\Leftrightarrow\; a = c \land b = d.
\]
\end{theorem}

\begin{proof}
($\Leftarrow$) Substitution.

($\Rightarrow$) Assume $\set{\set{a}, \set{a, b}} = \set{\set{c}, \set{c, d}}$.

\emph{Step 1: $a = c$.} $\set{a}$ is an element of the left side, so of the right
side: $\set{a} = \set{c}$ or $\set{a} = \set{c, d}$. In both cases $c \in
\set{a}$, so $c = a$.

\emph{Step 2, case $a = b$.} Then $(a, b) = \set{\set{a}, \set{a}} =
\set{\set{a}}$. The element $\set{c, d}$ of the right side must equal $\set{a}$,
so $d \in \set{a}$, so $d = a = b$.

\emph{Step 2, case $a \neq b$.} $\set{a, b}$ is an element of the right side
$\set{\set{a}, \set{a, d}}$ (using $c = a$). It is not $\set{a}$, since $b \in
\set{a, b}$ but $b \notin \set{a}$. So $\set{a, b} = \set{a, d}$. Then $b \in
\set{a, d}$ and $b \neq a$, so $b = d$.
\end{proof}

\begin{definition}[Tuples]
\label{def:tuple}
Let $(a, b, c) \coloneqq ((a, b), c)$, and generally
\[
  (a_1, \dots, a_{n+1}) \coloneqq ((a_1, \dots, a_n), a_{n+1}).
\]
Applying \cref{thm:pair} repeatedly, two tuples are equal if and only if their
components are equal one by one.
\register{relations}{$(a, b, c)$, \dots}{tuple $((a, b), c)$, and so on}
\end{definition}

\section{Cartesian products}

\begin{theorem}
\label{thm:product-exists}
For sets $A, B$ there is a unique set whose elements are exactly the pairs $(a,
b)$ with $a \in A$ and $b \in B$.
\end{theorem}

\begin{proof}
We need a set we already have that contains all such pairs, so that Separation
can pick them out. Let $a \in A$ and $b \in B$. Then $\set{a} \subseteq A \cup B$
and $\set{a, b} \subseteq A \cup B$, so both are elements of $\Pow (A \cup B)$.
Hence $(a, b) \subseteq \Pow (A \cup B)$, so $(a, b) \in \Pow \Pow (A \cup B)$. The
set
\[
  \setof{p \in \Pow \Pow (A \cup B)}{\exists a \in A\; \exists b \in B\; p = (a,
  b)}
\]
has the required elements. Uniqueness is by \cref{ax:extensionality}.
\end{proof}

\begin{definition}[Cartesian product]
\label{def:product}
The set from \cref{thm:product-exists} is the \emph{Cartesian product} $A \times
B$. For a predicate $P(x, y)$ we write
\[
  \setof{(x, y) \in A \times B}{P(x, y)} \coloneqq \setof{p \in A \times
  B}{\exists x \in A\; \exists y \in B\; (p = (x, y) \land P(x, y))}.
\]
\register{relations}{$A \times B$}{the set of pairs $(a, b)$ with $a \in A$, $b
\in B$}
\register{relations}{$\setof{(x, y) \in A \times B}{P(x, y)}$}{the set of pairs
in $A \times B$ satisfying $P$}
\end{definition}

\section{Relations}

\begin{definition}[Relation]
\label{def:relation}
A \emph{relation from $A$ to $B$} is a subset $R \subseteq A \times B$. A
\emph{relation on $A$} is a relation from $A$ to $A$. For $(x, y) \in R$ we write
\[
  x \mathrel{R} y,
\]
and by \cref{rule:abbreviation} the form $(x, y) \in R$ is not used again. The
\emph{domain}, \emph{range} and \emph{inverse} of $R$ are
\begin{gather*}
  \dom R \coloneqq \setof{x \in A}{\exists y \in B\; x \mathrel{R} y}, \qquad
  \ran R \coloneqq \setof{y \in B}{\exists x \in A\; x \mathrel{R} y}, \\
  R^{-1} \coloneqq \setof{(y, x) \in B \times A}{x \mathrel{R} y}.
\end{gather*}
So $R^{-1}$ is a relation from $B$ to $A$.
\register{relations}{$x \mathrel{R} y$}{$x$ is related to $y$ by $R$, i.e.\ $(x, y)
\in R$}
\register{relations}{$\dom R$, $\ran R$}{domain and range of $R$}
\register{relations}{$R^{-1}$}{inverse relation: $y \mathrel{R^{-1}} x
\Leftrightarrow x \mathrel{R} y$}
\end{definition}

\begin{definition}[Identity relation]
\label{def:identity}
The \emph{identity} on $A$ is $\id_A \coloneqq \setof{(x, y) \in A \times A}{x =
y}$.
\register{relations}{$\id_A$}{identity on $A$: $x \mathrel{\id_A} y
\Leftrightarrow x = y$}
\end{definition}

Other examples on a set $A$: the empty relation $\varnothing$ and the full
relation $A \times A$. On $\Pow A$ the statement $X \subseteq Y$ is itself a
relation, the set $\setof{(X, Y) \in \Pow A \times \Pow A}{X \subseteq Y}$, and
\cref{rule:relations} lets $\subseteq$ name both.

\begin{definition}[Properties of relations]
\label{def:relation-properties}
A relation $R$ on $A$ is
\begin{itemize}
  \item \emph{reflexive} if $\forall x \in A\; x \mathrel{R} x$;
  \item \emph{symmetric} if $\forall x, y \in A\; (x \mathrel{R} y \Rightarrow y
    \mathrel{R} x)$;
  \item \emph{antisymmetric} if $\forall x, y \in A\; (x \mathrel{R} y \land y
    \mathrel{R} x \Rightarrow x = y)$;
  \item \emph{transitive} if $\forall x, y, z \in A\; (x \mathrel{R} y \land y
    \mathrel{R} z \Rightarrow x \mathrel{R} z)$;
  \item \emph{irreflexive} if $\forall x \in A\; \lnot(x \mathrel{R} x)$.
\end{itemize}
\end{definition}

\section{Equivalence relations}

An equivalence relation says when two objects are \enquote{the same in the respect
we care about}. It is the main tool of the number constructions: integers and
rationals are both \emph{equivalence classes}.

\begin{definition}[Equivalence relation]
\label{def:equivalence}
An \emph{equivalence relation} on $A$ is a relation $\sim$ on $A$ that is
reflexive, symmetric and transitive. For $a \in A$, the \emph{equivalence class}
of $a$ is
\[
  [a] \coloneqq \setof{x \in A}{x \sim a}.
\]
The \emph{quotient set} is the set of all classes:
\[
  A/{\sim} \coloneqq \setof{C \in \Pow A}{\exists a \in A\; C = [a]}.
\]
Any $x \in C$ is a \emph{representative} of the class $C$.
\register{relations}{$[a]$}{equivalence class of $a$ (the relation is fixed by
the context)}
\register{relations}{$A/{\sim}$}{set of equivalence classes of $\sim$}
\end{definition}

\begin{lemma}
\label{lem:classes}
Let $\sim$ be an equivalence relation on $A$ and $a, b \in A$.
\begin{enumerate}
  \item $a \in [a]$.
  \item $[a] = [b] \;\Leftrightarrow\; a \sim b$.
  \item Either $[a] = [b]$ or $[a] \cap [b] = \varnothing$.
\end{enumerate}
\end{lemma}

\begin{proof}
(1) Reflexivity.

(2) ($\Rightarrow$) $a \in [a] = [b]$, so $a \sim b$. ($\Leftarrow$) Let $x \in
[a]$, so $x \sim a$. With $a \sim b$ and transitivity, $x \sim b$, so $x \in [b]$.
So $[a] \subseteq [b]$. By symmetry $b \sim a$, and the same argument gives $[b]
\subseteq [a]$. Use \cref{thm:subset}.

(3) Suppose some $x$ lies in both $[a]$ and $[b]$. Then $x \sim a$ and $x \sim b$,
so $a \sim x$ (symmetry) and $a \sim b$ (transitivity). By (2), $[a] = [b]$.
\end{proof}

So every element sits in exactly one class, and \emph{which} representative we
use to name a class doesn't matter.

\begin{definition}[Partition]
\label{def:partition}
A \emph{partition} of $A$ is a set $\Pi$ of nonempty subsets of $A$ such that
every $x \in A$ lies in exactly one $C \in \Pi$.
\end{definition}

\begin{theorem}[Equivalences are partitions]
\label{thm:partition}
\leavevmode
\begin{enumerate}
  \item If $\sim$ is an equivalence relation on $A$, then $A/{\sim}$ is a
    partition of $A$.
  \item If $\Pi$ is a partition of $A$, then the relation $\approx$ on $A$ given
    by $x \approx y \defiff \exists C \in \Pi\; (x \in C \land y \in C)$ is an
    equivalence relation whose classes are exactly the elements of $\Pi$.
\end{enumerate}
\end{theorem}

\begin{proof}
(1) Every class $[a]$ is nonempty, since it contains $a$
(\cref{lem:classes}(1)). Every $x \in A$ is in the class $[x]$, and if $x \in [a]$
and $x \in [b]$, then $[a] = [b]$ by \cref{lem:classes}(3).

(2) Reflexive: $x$ lies in some $C \in \Pi$. Symmetric: the definition is
symmetric in $x, y$. Transitive: if $x, y \in C$ and $y, z \in C'$, then $y$ is in
both $C$ and $C'$, so $C = C'$ by uniqueness, and $x, z \in C$. For the classes:
let $C$ be the element of $\Pi$ containing $a$. Then $x \approx a$ if and only if
$x$ and $a$ share an element of $\Pi$, if and only if $x \in C$ (uniqueness
again). So $[a] = C$, and every $C \in \Pi$ arises this way because $C \neq
\varnothing$.
\end{proof}

\begin{definition}[Well-defined]
\label{def:well-defined}
Often we define something on classes by using a representative: \enquote{for a
class $C$, pick $a \in C$ and let the result be \dots{} $a$ \dots}. This is a
valid definition only if the result is the same for \emph{every} choice of
representative. Proving that is called \emph{showing that the definition is
well-defined}.
\end{definition}

We will do this many times in \cref{ch:integers,ch:rationals}.

\section{Orders}

\begin{definition}[Partial and total order]
\label{def:order}
A relation $\le$ on $A$ is a \emph{partial order} if it is reflexive,
antisymmetric and transitive. Then $(A, \le)$ is a \emph{partially ordered set}.
It is a \emph{total order} if in addition any two elements are comparable:
\[
  \forall x, y \in A\; (x \le y \lor y \le x).
\]
We write $x < y$ for $x \le y \land x \neq y$. By \cref{rule:relations} there is
no separate symbol for the reversed relations.
\register{relations}{$x \le y$}{a partial order}
\register{relations}{$x < y$}{$x \le y$ and $x \neq y$}
\end{definition}

For example, $\subseteq$ on $\Pow A$ is a partial order by \cref{thm:subset}. It is
not total in general: on $\Pow \set{a, b}$ with $a \neq b$, neither $\set{a}
\subseteq \set{b}$ nor $\set{b} \subseteq \set{a}$.

\begin{lemma}
\label{lem:strict-order}
If $\le$ is a partial order, then $<$ is irreflexive and transitive. Conversely,
if $<$ is an irreflexive transitive relation on $A$, then $x \le y \defiff x < y
\lor x = y$ is a partial order.
\end{lemma}

\begin{proof}
Irreflexivity of $<$ is built into its definition. Transitivity: if $x < y$ and $y
< z$, then $x \le z$. If $x = z$, then $x \le y$ and $y \le x$, so $x = y$ by
antisymmetry, contradicting $x < y$. So $x < z$.

Conversely, $\le$ is reflexive by definition. Transitivity follows by cases on
which side of each \enquote{or} holds. Antisymmetry: if $x \le y$, $y \le x$ and
$x \neq y$, then $x < y$ and $y < x$, so $x < x$ by transitivity, contradicting
irreflexivity.
\end{proof}

\begin{definition}[Bounds]
\label{def:bounds}
Let $(A, \le)$ be partially ordered and $S \subseteq A$.
\begin{itemize}
  \item $m \in S$ is a \emph{least element} of $S$ if $m \le s$ for all $s \in S$,
    and a \emph{greatest element} if $s \le m$ for all $s \in S$.
  \item $u \in A$ is an \emph{upper bound} of $S$ if $s \le u$ for all $s \in S$,
    and a \emph{lower bound} if $u \le s$ for all $s \in S$. $S$ is \emph{bounded
    above} (\emph{below}) if it has an upper (lower) bound.
  \item A least element of the set of upper bounds of $S$ is a \emph{supremum}
    $\sup S$. A greatest element of the set of lower bounds is an \emph{infimum}
    $\inf S$.
\end{itemize}
\register{relations}{$\sup S$, $\inf S$}{least upper bound, greatest lower bound
of $S$}
\end{definition}

\begin{lemma}
\label{lem:least-unique}
A set has at most one least element. So $\sup S$ and $\inf S$ are unique when they
exist.
\end{lemma}

\begin{proof}
If $m$ and $m'$ are least elements of $S$, then $m \le m'$ and $m' \le m$, so $m
= m'$ by antisymmetry. The same works for greatest elements.
\end{proof}

Suprema become central with the real numbers: $\R$ is constructed so that
\emph{every} nonempty set bounded above has a supremum, which $\Q$ lacks
(\cref{thm:sqrt2}).

\begin{definition}[Well-order]
\label{def:well-order}
A total order $\le$ on $A$ is a \emph{well-order} if every nonempty subset of $A$
has a least element.
\end{definition}

\begin{exercises}
  \item Show that $(a, b, c) = (a', b', c')$ if and only if $a = a'$, $b = b'$ and
    $c = c'$.
  \item Show $A \times (B \cup C) = (A \times B) \cup (A \times C)$.
  \item On a set $A$, is the empty relation reflexive? Symmetric? Transitive?
    (Watch out when $A = \varnothing$.)
  \item Find an example showing that \enquote{symmetric and transitive} does not
    imply reflexive.
  \item In $(\Pow A, \subseteq)$, show that every $\mathcal{F} \subseteq \Pow A$
    has supremum $\bigcup \mathcal{F}$.
\end{exercises}
